\chapter{Methodology}

\section{Introduction}
This chapter describes the detection pipeline proposed in this thesis and the
systematic procedure used to build it. The design goal is a compact detector
that recovers very small human targets in aerial search-and-rescue imagery
without inflating the parameter or latency budget. Rather than adopting a single
fixed architecture, the model is developed through an \emph{additive ablation}:
starting from the YOLO11m baseline (reviewed in Chapter~II), one architectural
change is introduced at a time so that the contribution of each component can be
attributed independently. The chapter first presents the overall framework as a
structural anchor, then details each proposed modification, the state-space neck
variant that was explored, the ablation protocol, and a running example that
follows a single image through the pipeline. The baseline YOLO11 architecture and
the attention, multi-scale, and state-space concepts are covered in Chapter~II;
here the focus is strictly on the proposed modifications. Three further methods from
the extended phase of the work close the chapter: a tiny-aware change to the
label-assignment step, a frequency-gated evidence branch that was explored and
rejected, and the supervised procedure used to adapt the final model to real
aerial footage.

\section{Overall Framework}
The complete detection pipeline is shown in Figure~\ref{fig:framework}. An input
aerial image is letterbox-resized to $640 \times 640$ and passed through the
YOLO11m backbone, in which the terminal C2PSA attention block is replaced by a
Convolutional Block Attention Module (CBAM). The resulting multi-scale features
are aggregated by a PAN--FPN neck and decoded by a four-scale detection head that
adds a high-resolution P2 level (stride~4) to the standard P3--P5 levels. A final
non-maximum-suppression (NMS) stage produces the human bounding boxes. The two
proposed modifications, the CBAM attention block and the P2 detection scale,
are highlighted in the figure; the optional state-space neck variant
(Section~\ref{sec:mamba}) replaces selected neck blocks and is evaluated
separately.

\begin{figure}[tbp]
  \centering
  \figorplaceholder{figures/fig_overall_framework.png}{6cm}
  \caption{Overall framework of the proposed detector. An aerial image is
  preprocessed and passed through the YOLO11m backbone (with CBAM replacing the
  C2PSA block after the SPPF stage), a PAN--FPN neck, and a four-scale detection
  head; the added P2 level at stride~4 targets very small targets, and NMS yields
  the final human detections.}
  \label{fig:framework}
\end{figure}

\section{Detailed Methodology}

\subsection{Data Preprocessing}
All images are letterbox-resized to $640 \times 640$ so that the original aspect
ratio is preserved and no target is geometrically distorted, which matters when
the objects of interest are only a few pixels wide. Pixel values are normalised to
$[0,1]$, and only the default YOLO11 training-time augmentations are used
(mosaic, scaling, HSV jitter), with mosaic disabled for the final ten epochs
(\texttt{close\_mosaic}). Aggressive small-object augmentation (copy--paste,
mixup, vertical flip) was tested and discarded as a negative result, as reported
in Chapter~IV.

\subsection{Attention Enhancement: CBAM}
The baseline YOLO11m places two C2PSA self-attention blocks (1024 channels) after
the SPPF stage. In the proposed model these are replaced by a single CBAM module,
which applies channel attention followed by spatial attention in sequence. Given
an intermediate feature map $\mathbf{F} \in \mathbb{R}^{C \times H \times W}$, CBAM
computes a channel attention map $\mathbf{M}_c$ and a spatial attention map
$\mathbf{M}_s$,
\begin{align}
  \mathbf{M}_c(\mathbf{F}) &= \sigma\!\big(\mathrm{MLP}(\mathrm{AvgPool}(\mathbf{F}))
      + \mathrm{MLP}(\mathrm{MaxPool}(\mathbf{F}))\big), \\
  \mathbf{M}_s(\mathbf{F}') &= \sigma\!\big(f^{7\times7}\big([\mathrm{AvgPool}(\mathbf{F}');
      \mathrm{MaxPool}(\mathbf{F}')]\big)\big),
\end{align}
where $\sigma$ is the sigmoid function and $f^{7\times7}$ a convolution with a
$7\times7$ kernel. The features are refined multiplicatively,
$\mathbf{F}' = \mathbf{M}_c(\mathbf{F}) \otimes \mathbf{F}$ and
$\mathbf{F}'' = \mathbf{M}_s(\mathbf{F}') \otimes \mathbf{F}'$, where $\otimes$
denotes element-wise multiplication with broadcasting. Because CBAM is lighter
than the two C2PSA blocks it replaces, this substitution slightly \emph{reduces}
the parameter count relative to the baseline while adding both channel and spatial
selectivity.

\subsection{Multi-Scale Detection: the P2 Head}
The standard YOLO11 head detects at three scales: P3, P4 and P5 (strides 8, 16
and 32). For $640 \times 640$ input, the finest of these (P3) operates on an
$80 \times 80$ grid, so each cell already summarises an $8 \times 8$ pixel region;
targets smaller than this are difficult to localise. The proposed model adds a
fourth detection level, P2, at stride~4. This introduces a $160 \times 160$
feature map in which each cell covers only a $4 \times 4$ pixel region, allowing
the detector to resolve targets as small as approximately four pixels. The P2
level is formed by up-sampling and fusing the corresponding high-resolution
backbone features into the neck and attaching an additional detection branch, as
summarised in Figure~\ref{fig:framework}. As shown in Chapter~IV, this
high-resolution scale is the principal driver of the improvement in very-small-%
target recall.

\subsection{State-Space Neck Variant: C3k2Mamba}
\label{sec:mamba}
To test whether a selective state-space model (SSM) can improve feature
aggregation in the neck, an explored variant replaces the C3k2 bottleneck at six
neck layers (indices 13, 16, 19, 22, 25 and 28) with a C3k2Mamba block. Each such
block performs bidirectional local-window scanning, a forward and a reverse
selective scan over a local window (window sizes 6--8, state dimension
$d_{\text{state}} = 4$), and is inserted by post-initialisation surgical
replacement so that the surrounding CNN backbone and the P2 head are unchanged.
This variant is included for completeness and evaluated under the same protocol;
Chapter~IV reports that it does not improve accuracy while increasing both the
parameter count and the inference latency, and it is therefore excluded from the
recommended model.

\subsection{The Additive Ablation Design}
The four configurations evaluated in this thesis are summarised in
Table~\ref{tab:variants}. Each row introduces exactly one change with respect to
the preceding design, so that the effect of attention (CBAM), of the
high-resolution detection scale (P2), of their combination, and of the
state-space neck can each be isolated. All configurations are trained on the same
frozen C2A split with identical optimisation settings (AdamW, $\text{lr}_0 = 0.001$,
cosine schedule, up to 300 epochs with early stopping) and evaluated with the same
protocol and thresholds, so that every reported difference is attributable to the
architectural change alone. The experimental setup and these settings are detailed
in Chapter~IV.

\begin{table}[tbp]
  \centering
  \caption{The additive ablation. Each configuration differs from the previous one
  by a single architectural change, isolating the contribution of attention,
  multi-scale detection, and the state-space neck.}
  \label{tab:variants}
  \begin{tabular}{@{}l l l l@{}}
    \toprule
    Configuration & Backbone attention & Detection scales & Neck \\
    \midrule
    YOLO11m (baseline)   & C2PSA            & P3, P4, P5         & C3k2 \\
    + CBAM               & CBAM             & P3, P4, P5         & C3k2 \\
    + CBAM + P2          & CBAM             & P2, P3, P4, P5     & C3k2 \\
    + Mamba + CBAM + P2  & CBAM             & P2, P3, P4, P5     & C3k2Mamba (6 layers) \\
    \bottomrule
  \end{tabular}
\end{table}

\subsection{Running Example}
It helps to trace one image through the pipeline of Figure~\ref{fig:framework}.
Consider a collapsed-building scene of the kind shown in
Figure~\ref{fig:dataset_samples}(a), a $2000\times1500$ aerial photograph with more
than thirty people scattered across rubble, most of them under sixteen pixels tall.
The image is first letterbox-resized to $640\times640$, which preserves the aspect
ratio so the already small figures are not stretched. The backbone extracts features
at five scales, and the CBAM block after the SPPF stage re-weights them so that
human-shaped regions are emphasised over the surrounding debris. The neck fuses the
scales, and the head predicts boxes at all four levels; the added P2 level, at
stride four, is the one that fires on the smallest figures, which the stride-eight
P3 level would summarise into a single cell. Non-maximum suppression then removes
overlapping boxes, so that a cluster of nearby people is kept as separate detections
rather than merged. The output is one bounding box per recovered person, which is
what an operator reviews.

\subsection{Tiny-Aware Label Assignment}
\label{sec:assignmethod}
The methods so far change the network. The extended study of
Section~\ref{sec:partb} adds three methods that leave the recommended
architecture untouched, described in this and the next two subsections. The
first acts on the step that decides which predictions are trained at all. During
training, the task-aligned assigner scores every candidate prediction against
every ground-truth box and selects the top candidates as positive samples, using
the alignment score
\begin{align}
  t = s^{\,p} \, u^{\,q},
\end{align}
where $s$ is the classification score of the candidate, $u$ is the localisation
overlap between the candidate box and the ground-truth box, and $p$ and $q$ are
fixed exponents. The default overlap is the complete IoU, which is the weak
point for this task: for a person a few pixels wide, a one-pixel offset between
candidate and ground truth can drive the IoU to zero, so the ranking among
candidates for the smallest targets carries almost no signal, and the smallest
people are trained on poorly chosen samples. Prior work on tiny-object
benchmarks reached the same diagnosis and moved the fix into the assignment
step, as reviewed in Chapter~II.

The modification keeps the assigner and replaces only its overlap term. For a
ground-truth box $g$ and a candidate $b$ with centres $(c_x, c_y)$, widths $w$,
and heights $h$, the squared 2-Wasserstein distance between the Gaussians fitted
to the two boxes is
\begin{align}
  W_2^2(g,b) = (c_{x}^{g}-c_{x}^{b})^2 + (c_{y}^{g}-c_{y}^{b})^2
             + \tfrac{1}{4}\big[(w^{g}-w^{b})^2 + (h^{g}-h^{b})^2\big],
\end{align}
and the normalised Wasserstein similarity~\cite{wang2021nwd} maps it into
$(0,1]$ as
\begin{align}
  S_{\mathrm{NWD}}(g,b) = \exp\!\big(-\,W_2(g,b)\,/\,C\big),
\end{align}
with the scale constant $C=12.8$ set to the median person size of the dataset in
pixels. Unlike the IoU, this similarity stays informative for disjoint tiny
boxes: a near-miss candidate scores higher than a distant one even when both
have zero overlap. The overlap term used by the assigner becomes the blend
\begin{align}
  u = (1-\alpha)\,\mathrm{CIoU}(g,b) + \alpha\, S_{\mathrm{NWD}}(g,b),
  \label{eq:blend}
\end{align}
where $\alpha \in [0,1]$ sets how far the assignment trusts the scale-robust
similarity over the IoU; $\alpha=0$ recovers the unmodified assigner.
Figure~\ref{fig:assignmech} places the blend inside the assignment pipeline. The
change affects training only, adds no parameters and no inference cost, and is
distinct from using the Wasserstein distance inside the regression loss, which
was also tried and did not help; Chapter~IV reports the sweep over $\alpha$ and
the matched comparison at $\alpha=0.5$.

\begin{figure}[tbp]
  \centering
  \figorplaceholder{figures/fig_assign_mechanism.png}{6.5cm}
  \caption{Tiny-aware label assignment. For each ground-truth and candidate box
  pair the assigner computes a classification score $s$ and a localisation
  overlap $u$; the proposed change replaces $u$ with a blend of the complete IoU
  and the normalised Wasserstein similarity, weighted by $\alpha$, before the
  usual top-$k$ selection of positive training samples. The network, the loss,
  and inference are unchanged.}
  \label{fig:assignmech}
\end{figure}

\subsection{A Frequency-Gated Evidence Branch (Explored)}
\label{sec:fccgmethod}
The second added method is an architecture branch that was built, tested, and
rejected; it is described so that the negative result of Chapter~IV is
interpretable. A degradation probe showed that the detector's tiny-object signal
lives almost entirely in the highest spatial-frequency band of the image, which
suggested making that band explicit. The branch extracts high-frequency
evidence from the finest feature maps and multiplies it by a gate computed with
large-kernel convolutions from the surrounding context, the intent being to keep
high-frequency responses where context indicates a plausible person and suppress
them over clutter. The gated evidence is added back at the two finest fusion
points of the neck (the P3 and P2 seams), adding $0.55$~M parameters. The gates
train stably and remain active, so the branch fails for a substantive reason
rather than a mechanical one: whatever it contributes, the CBAM~+~P2 features
already carry.

\subsection{Supervised Adaptation to Real Imagery}
\label{sec:adaptmethod}
The third added method carries the trained detector from composited training
data to real footage. Real aerial frames were collected from three sources: the
author's own drone flights over open ground at 10, 30, and 50~m altitude,
recorded in 4K; campus scenes with partial occlusion from the same platform; and
frames sampled from published news footage of two disaster events, a flood and a
building collapse. Frames were extracted at fixed intervals, de-duplicated with
a perceptual hash, and annotated by the author with segmentation-assisted box
drawing in a web labelling tool, with every visible person boxed and frames
containing no people kept as explicit negatives. An automated audit
(degenerate, duplicate, undersized, and oversized boxes) and a visual overlay
check of every frame preceded any training use. Evaluation frames are frozen:
the same perceptual hash guarantees that no training frame is a near-duplicate
of any evaluation frame.

Adaptation is a supervised fine-tune of the assignment model from
Section~\ref{sec:assignmethod}. The training mixture holds the full C2A training
split together with the real training frames repeated twenty times, so that the
few hundred real frames are not drowned out by the 6{,}135 synthetic images while
the synthetic distribution is retained; 4K frames are downscaled to 1{,}280~px on
the long side beforehand. The fine-tune runs at a reduced learning rate
($\mathrm{lr}_0 = 5\times10^{-4}$, cosine schedule) for at most 80 epochs with
the same dual early stopping as the main protocol, and the assignment blend of
Equation~\eqref{eq:blend} stays active. This recipe deliberately contains no
domain-adaptation machinery beyond data mixing: it measures what a small
labelled budget buys with the simplest reliable procedure, which makes the
result a floor for more elaborate methods rather than a claim of novelty in the
adaptation itself.

\section{Conclusion}
This chapter presented the proposed detector as an additive refinement of YOLO11m:
CBAM attention in place of C2PSA, a high-resolution P2 detection scale for very
small targets, and an explored state-space neck variant. The overall framework
anchors these modifications, and the additive ablation
design ensures that each component's contribution can be measured independently.
The extended methods change how the detector is trained rather than what it is:
a Wasserstein-based blend inside the label assigner, a rejected frequency-gated
branch, and a data-mixing fine-tune that adapts the model to real footage. The
next chapter reports the experimental setup and the quantitative and
qualitative results of the ablation and of the extended study.
